\section{Quadratic original models after the Type II exclusion}
\label{sec:cubic-quadratic-model-rigidity}

The exclusion is about actual cubic Type II points as well as about
possible rational starting points.  This distinction removes a tangent
premise from the original quadratic model once its actual Type I
support is supplied.  The theorem below specifies that support directly;
it does not retain an impossible original quadratic seed as a hypothesis.

Put $p=269$, $K=\mathbb Q(\sqrt{-3})$, and let $\chi$ be either
nontrivial character of $K(\sqrt[3]{269})/K$.  Use the original
coefficient DVR $O$, the actual normalized cubic ring $S_j$, and its
graded model $(A,\mathcal A,x,q_A,\mathcal Q)$ from AC-G65--AC-G66.
Write $I=\mathcal Q.\mathrm{ideal}$ and let $L_j$ be the actual
weight-minus-one cotangent length.  Every original local, marked-element,
cohomological and model hypothesis is retained.  In particular the
coefficient condition $e_O=1$ required by the old modularity input is
imposed on this original $O$, not on a later value or host coefficient
ring.  Retain the published local/old modularity inputs in AC-G117,
including the literal $\mathrm{GD}$ and
\texttt{AutomorphyUnramifiedCoefficientsAlt}.

For the first theorem take any original quadratic character
$\psi^2=1$, including the trivial character, and any weight admitted
by that datum.  Assume the actual model condition $\mathrm{M2}$ and
an actual support host and native modular ideal satisfying exactly
(G110.2), with every native classicality and published adjoint input of
AC-G110.  These are statements about the actual cubic model and its
actual host maps.  Merely giving a closed automorphic point, an abstract
finite algebra, or a rational auxiliary form does not establish them.

\begin{theorem}[AC-G122-M4: rigidity after quadratic Type II exclusion]
\label{AC-G122-M4}
Under these hypotheses every actual two-sided characteristic-zero
local point of $S_j$ has Type I and is rigid.  The original model
satisfies the literal source condition
\[
                \mathrm{M4}(\mathcal A,\overline{\mathbb Q}_p,I).
                                                        \tag{G122.1}
\]
No mixed Type II Selmer-vanishing hypothesis or
\texttt{RigiditySecondCase} hypothesis is needed for this conclusion.
The assertion quantifies over all finite value-field extensions allowed
by the actual model.  It neither changes $O$ nor asserts vanishing for
a mixed coefficient attached to a point outside this excluded family.
\end{theorem}

\begin{proof}
\proofstep{1. Exclude the actual second case.}
The normalized shape argument of AC-G113, Step~1 of its M5 proof,
uses the actual diagonal difference
$\operatorname{diag}(\zeta_f,0,-\zeta_f)$ with $\zeta_f\ne0$.
Its four possibilities are Type I, upper, lower and Type II.
The actual model signs exclude upper and lower at a two-sided point.
If Type II remained, AC-G68-block and AC-G68-point would untwist its
endpoint block by $\chi^{-j}$ and construct an actual original rational
endpoint with every original inertia, ordinary, selected-ideal and
modularity property.  AC-G117 excludes precisely that endpoint for
every original quadratic $\psi$ under the retained hypotheses.
The contradiction is unchanged by enlarging the value field: the
original coefficient datum in AC-G68-point stays $O$, and the
exclusion only uses the resulting actual representation and its
residual, local and modular properties.  Thus every such point is
Type I.

AC-G110-rigidity now applies to each point using the actual native
support (G110.2).  Its original rational adjoint Bloch--Kato group
vanishes, and AC-G66-BK puts every actual dual-point tangent in that
group.  Consequently every actual dual-point cocycle is a coboundary.
This proves the asserted point rigidity without constructing any
mixed Type II group or invoking its vanishing.

\proofstep{2. Check the full derivation quantifier in M4.}
Let $\xi:A\to\overline{\mathbb Q}_p$ kill $I$ and be two-sided.
Let $D,D'$ be two $O$-linear point derivations at $\xi$, both
annihilating $I$.  The actual model point property, using M2, gives
$t\in\overline{\mathbb Q}_p^\times$ and a local point
$f:S_j\to\overline{\mathbb Q}_p$ such that
\[
                   \operatorname{act}(t)\xi=fq_A.
                                                        \tag{G122.2}
\]
For a homogeneous element $a\in\mathcal A_d$ set
\[
                 E(a)=t^dD(a),\qquad E'(a)=t^dD'(a),
                                                        \tag{G122.3}
\]
and extend by the direct-sum grading.  The identity
$t^{d+d'}=t^dt^{d'}$ and the product rule show that these are point
derivations at $fq_A$.  Every specified generator of $I$ is
homogeneous.  Its derivative is zero by (G122.3); the product rule
and $fq_A(I)=0$ therefore show that both derivations kill the entire
ideal $I$.  The actual model dual-point extension property gives
dual maps on $S_j$ inducing $E,E'$.

The factors $t^d$ are nonzero, so both sign conditions are preserved
in (G122.2).  The source model-sign argument, with the actual entry
grading, makes $f$ a two-sided point.  Step~1 makes its two dual
cocycles coboundaries.  The normalization at complex conjugation and
the marked difference then invokes precisely
\texttt{normCoc\_cob\_iff} in \texttt{NormCob.lean}.
Its hypotheses are that $2$ and $\zeta_f$ are units, the multiplier
is $-1$ at complex conjugation and $1$ at the marked element, and
the representation is normalized.  All hold here.  Its conclusion
puts each cobounding matrix in
\[
        \overline{\mathbb Q}_p\operatorname{diag}(1,0,-1).
                                                        \tag{G122.4}
\]
The two cocycles are therefore linearly dependent, with coefficients
$a,a'$ not both zero.  The actual dual-entry determination property
of AC-G66, equivalently the applied \texttt{tanCoc\_determines}
step in the source derivation proof, makes the same linear
combination of dual derivatives zero on $S_j$.  Restricting along
$q_A$ gives $aE+a'E'=0$.  On degree $d$, cancel $t^d$ in
(G122.3); the direct-sum grading then gives
$aD+a'D'=0$ on all of $A$.

This is every quantifier in the definition of M4 in
\texttt{OrbitVanishingLocal.lean}.  The argument is the complete
generic derivation mechanism in
\texttt{GradedModelDerivationCondition\allowbreak OfPtHyp.lean}, applied to the
actual model properties, rather than declaring its old
single-character \texttt{PtHyp} for the cubic ring.  It proves
(G122.1).
\end{proof}

\paragraph{The concrete branch and its retained support boundary.}
For the remainder use
\[
        N=4,\quad\psi=\chi_{-4},\quad
        e\equiv267\pmod{268},\quad e\ge1,
        \quad j\in\{1,2\}.
                                                        \tag{G122.5}
\]
Keep the actual M2 and Type I support just specified, the remaining
structural hypotheses of AC-G113, and the rank-three descent,
local-compatibility, witness and counting inputs of AC-G111 and
AC-G112-count.  In particular the support host has a CM tower with
cyclic steps of degree $2$ or $269$, as required by the exact
descent theorem.  Retain the generic algebraic statements
\texttt{LemmaM} and \texttt{LemmaM'}.
No original quadratic rational seed is assumed here.  A construction
using the auxiliary seed must prove these actual support hypotheses
separately; the present theorem does not silently infer them from G119.
The old rank-two trace classifier and its counting hypotheses are
unnecessary below, since its possible actual points have been excluded.

\begin{theorem}[AC-G122-count: the Type I integral model bound]
\label{AC-G122-count}
The original model in (G122.5) satisfies M1--M5.  Its cotangent
length is finite.  With $\sigma$ geometric Frobenius at the unique
place of $K$ over $11$, set
\[
                   b=121^e+1+121^{-e},\qquad Y=4(e+1).
                                                        \tag{G122.6}
\]
There is a finite family $X_j$ of actual Type I representations,
pairwise nonisomorphic on $G_K$, with
\[
 s_j\le\frac{L_j}{e_O}
       \le2\sum_{\rho\in X_j}
                  v_{269}(\operatorname{tr}\rho(\sigma)-b),
       \qquad\#X_j\le C_1Y^{A_1}.
                                                        \tag{G122.7}
\]
Here $s_j$ is the exact original strict-group normalized length
in (G113.11), and $C_1,A_1$ are the rank-three constants of
AC-G112-count.  For every nonempty $J_0\subset\{1,2\}$ to which
the displayed hypotheses apply, there is a common nonzero integer
$\mathcal N^I_{11,J_0}$ with
\begin{align}
 \sum_{j\in J_0}\frac{L_j}{e_O}
    &\le2v_{269}(\mathcal N^I_{11,J_0}),\notag\\
 \log|\mathcal N^I_{11,J_0}|
    &\le5C_1\#J_0Y^{A_1+1}\log11,\notag\\
 \left(\sum_{j\in J_0}\frac{L_j}{e_O}\right)\log269
    &\le10C_1\#J_0Y^{A_1+1}\log11.
                                                        \tag{G122.8}
\end{align}
The same upper bounds hold with $L_j/e_O$ replaced by $s_j$.
The trace classifier and norm integer are selected before the
coefficient prime; the model application here remains at $269$.
\end{theorem}

\begin{proof}
AC-G66 supplies M1 and M3; M2 is the actual retained support
consequence; AC-G122-M4 supplies M4.  Every two-sided local point
is Type I.  The Type I part of AC-G113-M5 consequently proves M5:
the actual degree-zero trace difference has value zero at the
origin, while its value at every two-sided point is nonzero
because $b>3$ and the actual normalized trace has complex absolute
value at most $3$.  The all-model-point passage uses (G122.2),
which acts trivially in degree zero.  No Type II classifier or
Type II purity assertion enters this argument.

Apply \texttt{LemmaM} with the exact valuation $v_O=e_Ov_{269}$.
The full first three steps of AC-G113-bound then give a finite
orbit-reduced cover and the comparison (G122.7), including the
strict integral injection from \texttt{LemmaM'}.  Their stronger
isomorphism comparison is the normalized diagonal-difference
argument and is valid for these actual Type I points.  Every
representation in the resulting $X_j$ has Type I by the first
theorem.  Hence the single rank-three classifier of AC-G112-count
gives $\#X_j\le C_1Y^{A_1}$; no disjoint union with a rank-two
carrier is formed.

Apply the exact source
\texttt{frobenius\_trace\_norm\_nonzero\_log\_bound} to this
\texttt{Af.TrData}, with norm $Q=121$, label set $J_0$, and
$z(j)=1$ for every label.  This label family is Galois stable.
The actual seed trace at this inert test place is exactly $b$,
as proved in AC-G113, Step~1: the cubic character is trivial
on Frobenius at this degree-two place.  The theorem gives a
nonzero integer selected before the coefficient prime with
\[
 |\mathcal N^I_{11,J_0}|_{269}
   \le\prod_{j\in J_0}\prod_{\rho\in X_j}
                  |\operatorname{tr}\rho(\sigma)-b|_{269}.
                                                        \tag{G122.9}
\]
Its logarithmic bound is
$\#J_0\#\mathcal T(2e\log121+1)$.
Since $\#\mathcal T\le C_1Y^{A_1}$, $e+1\le Y$ and
$\log11>1$, this is at most the second line of (G122.8).
Taking negative logarithms in (G122.9), and summing the middle
inequality of (G122.7), gives the first line of (G122.8), for
the model lengths themselves.  Finally
$v_{269}(n)\log269\le\log|n|$ for every nonzero rational integer
$n$ gives the third line.  This proves all claims.
\end{proof}

\paragraph{Arithmetic boundary and geometric dictionary.}
AC-G114 already gives $s_1=s_2=0$ directly on (G122.5), without
model or support premises.  The new conclusion here concerns the
actual model tangent condition and its cotangent lengths.  It
does not turn that single branch into the all-prime, all-character,
all-weight bound H04.4.  On other families an actual Type II
point may exist and its mixed tangent coefficient remains a
separate question.

Geometrically M4 says that, at every two-sided characteristic-zero
point of this actual graded affine model, the tangent space cut
out by its equations has dimension at most one.  The surviving
direction is the infinitesimal diagonal-torus action.  The finite
cover and local-system identifications used to obtain the point
and its adjoint are those of AC-G68 and AC-G110, in the already
audited SGA~1, Expose~V, sections~4 and~7.  The orbit argument is
an assertion on this deformation chart, not a new arithmetic
square, motivic quotient or general cohomological Hodge theorem.

\paragraph{Exact proof and source scope.}
The complete proofs of AC-G113, AC-G110, the applied AC-G68
interface and AC-G112-count have been reread here; the whole
AC-G117 review is retained.  Fresh complete containing reads of
\texttt{OrbitVanishingLocal.lean},
\texttt{GradedModelDerivationCondition\allowbreak OfPtHyp.lean},
\texttt{NormCob.lean}, \texttt{ModelUpper.lean}, and the exact
\texttt{Source\_FrobNorm.lean} theorem at the pinned source
commit check the stated quantifiers.  The published modularity,
classicality, local comparison, cyclic descent, adjoint and
counting inputs remain at their explicit prior scopes.  No Lean
build or new proof of those upstream theorems is claimed.
